\documentclass[../../../include/open-logic-section]{subfiles}

\begin{document}

\olfileid[id]{sth}{replacement}{limofsize}
\olsection{Pembatasan Ukuran}

Mungkin upaya yang paling umum untuk memberikan pembenaran ``intrinsik'' bagi
Penggantian menggunakan gagasan berikut:
\begin{enumerate}
	\item[] \limofsize. Sembarang hal membentuk suatu himpunan, asalkan
	jumlahnya tidak terlalu banyak.
\end{enumerate}
Prinsip ini akan langsung membenarkan Penggantian. Bagaimanapun, setiap
himpunan yang dibentuk melalui Penggantian tidak mungkin lebih besar daripada
himpunan asal pembentukannya. Dinyatakan secara tepat: andaikan Anda
membentuk suatu himpunan
$\funimage{\tau}{A} = \Setabs{\tau(x)}{x \in A}$ dengan menggunakan
Penggantian; maka $\cardle{\funimage{\tau}{A}}{A}$; jadi, jika !!{element}s
dari $A$ tidak terlalu banyak untuk membentuk suatu himpunan, citranya juga
tidak terlalu banyak untuk membentuk $\funimage{\tau}{A}$.

Kesulitan yang jelas dalam menggunakan \limofsize{} untuk membenarkan
Penggantian ialah bahwa kita \emph{belum} menetapkan prinsip apa pun seperti
\limofsize. Selain itu, ketika kita menceritakan kisah tentang konsepsi
kumulatif-iteratif mengenai himpunan dalam
\crefrange{sth:story::chap}{sth:z::chap}, tidak ada apa pun yang pernah
\emph{mengisyaratkan} \limofsize. Inilah persis alasan Boolos pernah menulis:
``Mungkin orang dapat menyimpulkan bahwa terdapat setidaknya dua gagasan
`di balik' teori himpunan'' \citeyearpar[p.~19]{Boolos1989}. Di satu pihak,
gagasan-gagasan seputar konsepsi kumulatif-iteratif tentang himpunan
dimaksudkan untuk membenarkan~$\Z$. Di pihak lain, \limofsize{} dimaksudkan
untuk membenarkan Penggantian.

Namun, masalahnya bukan sekadar bahwa sejauh ini kita \emph{diam} mengenai
\limofsize. Alih-alih, masalahnya ialah bahwa \limofsize{} (sebagaimana baru
saja dirumuskan) tampaknya sangat tidak selaras dengan gagasan
kumulatif-iteratif tentang himpunan. Bagaimanapun, prinsip itu sama sekali
tidak menyebut gagasan bahwa himpunan dibentuk dalam \emph{tahap-tahap}.

Hal ini sebenarnya tidak terlalu mengejutkan, mengingat sejarah kedua
``gagasan'' ini (yakni konsepsi kumulatif-iteratif tentang himpunan dan
\limofsize). Kedua ``gagasan'' ini pada akhirnya merupakan dua proyek yang
cukup berbeda untuk menghalangi paradoks-paradoks teori himpunan. Gagasan
kumulatif-iteratif tentang himpunan menghalangi Paradoks Russell dengan
mengatakan, secara garis besar: \emph{kita tidak semestinya pernah
mengharapkan himpunan Russell ada, karena himpunan itu tidak akan
``terbentuk'' pada tahap mana pun}. Sebaliknya, \limofsize{} dimaksudkan
untuk menyingkirkan himpunan Russell dengan mengatakan, secara garis besar:
\emph{kita tidak semestinya pernah mengharapkan himpunan Russell ada, karena
himpunan itu terlalu besar}.

Jika diungkapkan seperti ini, mari kita berterus terang: sebagai jawaban
terhadap paradoks-paradoks, \limofsize{} memerlukan \emph{jauh} lebih banyak
pembenaran. Sebagai contoh, perhatikan versi Paradoks Russell berikut:
\emph{tidak ada anjing pug yang mengendus tepat anjing-anjing pug yang tidak
mengendus dirinya sendiri} (lihat \olref[sth][story][rus]{sec}). Jika Anda
bertanya ``mengapa tidak ada anjing pug semacam itu?'', diberi tahu bahwa
anjing tersebut harus mengendus terlalu banyak anjing pug bukanlah jawaban
yang baik. Jadi, mengapa ketiadaan himpunan Russell akan dijelaskan secara
intuitif dengan baik oleh klaim bahwa himpunan itu akan terlalu ``besar''
untuk ada?

Singkatnya, dapat dimaklumi jika Anda merasa agak bingung mengenai motivasi
``intuitif'' bagi \limofsize.

\end{document}
